\documentclass[11pt]{article}
\usepackage[OT1]{fontenc}
\usepackage[utf8]{inputenc}
\usepackage{amsmath,amssymb,amsthm,booktabs,geometry,listings,xcolor,graphicx,needspace,tikz,longtable}
\usepackage[hidelinks]{hyperref}
\usetikzlibrary{arrows.meta,calc,decorations.pathreplacing,fit,matrix,positioning,shapes.geometric}
\geometry{margin=1in}
\newtheorem{theorem}{Theorem}
\newtheorem{lemma}[theorem]{Lemma}
\newtheorem{definition}[theorem]{Definition}
\newcommand{\F}{\mathbb F_2}
\newcommand{\wt}{\operatorname{wt}}
\newcommand{\spanof}{\operatorname{span}}
\newcommand{\cost}{\ell}
\definecolor{proofnavy}{HTML}{18324B}
\definecolor{proofblue}{HTML}{2468A2}
\definecolor{proofpaleblue}{HTML}{EDF4FA}
\definecolor{proofgreen}{HTML}{087A68}
\definecolor{proofpalegreen}{HTML}{EDF8F4}
\definecolor{prooforange}{HTML}{B45A08}
\definecolor{proofpaleorange}{HTML}{FFF4E8}
\definecolor{proofgray}{HTML}{667482}
\definecolor{proofpalegray}{HTML}{F3F5F7}
\definecolor{proofgrid}{HTML}{B8C4CF}
\tikzset{
  proof title/.style={font=\bfseries\large,text=proofnavy,align=center},
  proof box/.style={draw=proofblue,fill=proofpaleblue,rounded corners=2pt,
    line width=.8pt,align=center,inner sep=5pt,text=proofnavy},
  proof green/.style={proof box,draw=proofgreen,fill=proofpalegreen},
  proof orange/.style={proof box,draw=prooforange,fill=proofpaleorange},
  proof gray/.style={proof box,draw=proofgray,fill=proofpalegray},
  proof arrow/.style={-{Stealth[length=2.2mm]},line width=.8pt,draw=proofgray},
  proof green arrow/.style={proof arrow,draw=proofgreen},
  proof orange arrow/.style={proof arrow,draw=prooforange},
  proof note/.style={font=\small,text=proofgray,align=center}
}
\setcounter{MaxMatrixCols}{13}
\newcommand{\prooffigure}[3]{%
  \begin{center}
  \begin{minipage}{.98\linewidth}
  \centering
  \resizebox{\linewidth}{!}{\input{figures/#1.tex}}\par
  \refstepcounter{figure}\label{#2}%
  \small\textbf{Figure \thefigure.} #3
  \end{minipage}
  \end{center}}
\lstset{basicstyle=\ttfamily\scriptsize,breaklines=true,columns=fullflexible,
  keepspaces=true,showstringspaces=false,numbers=left,numberstyle=\tiny}
\title{Proving the Exact Minimum CNOT Length\\of a $13\times13$ Binary Transformation}
\author{Ivan Litvak\\Independent Researcher\\
\texttt{vanya.litvak@gmail.com}}
\date{27 September 2026}
\begin{document}
\maketitle
\begin{abstract}
We establish that the minimum CNOT length of the explicit binary matrix
below is nineteen. Structural lower bounds from projected XOR complexity,
forced copies and restorations, and clean-carrier invariants reduce the
universal lower bound to 363 finite cases. Exact passive-target elimination
then permits an exhaustive external C++ check; a nineteen-gate witness
establishes attainability. Lean formalizes the general reduction and proves
the minimum conditional on two precisely stated finite facts about the
concrete lookup tables and search. We include the verifier, complete-run
evidence, and Lean sources. The finite facts are supported by an audited
external computation, not replayed by the Lean kernel.
\end{abstract}

\section{The matrix and the exact-minimum statement}
Let $\F=\{0,1\}$ with addition modulo two. The input is an arbitrary vector
$x=(x_0,\ldots,x_{12})^\top\in\F^{13}$; its thirteen coordinates are
independent variables, not values restricted to a code space. We study the
full labelled linear map $x\mapsto Mx$. With both rows and columns ordered
$0,1,\ldots,12$, the matrix is explicitly
\begin{equation}\label{eq:matrix}
M=\begingroup\setlength{\arraycolsep}{4pt}\renewcommand{\arraystretch}{.85}
\begin{pmatrix}
1 & 0 & 0 & 0 & 0 & 0 & 0 & 0 & 0 & 0 & 0 & 0 & 0 \\
0 & 1 & 0 & 0 & 0 & 0 & 0 & 0 & 0 & 0 & 0 & 0 & 0 \\
0 & 0 & 1 & 0 & 0 & 0 & 0 & 0 & 0 & 0 & 0 & 0 & 0 \\
0 & 0 & 0 & 1 & 0 & 0 & 0 & 0 & 0 & 0 & 0 & 0 & 0 \\
0 & 0 & 0 & 0 & 1 & 0 & 0 & 0 & 0 & 0 & 0 & 0 & 0 \\
0 & 0 & 0 & 0 & 0 & 1 & 0 & 0 & 0 & 0 & 0 & 0 & 0 \\
1 & 0 & 0 & 1 & 0 & 0 & 1 & 0 & 1 & 0 & 0 & 0 & 0 \\
0 & 1 & 0 & 0 & 1 & 0 & 0 & 1 & 1 & 0 & 0 & 0 & 0 \\
0 & 0 & 1 & 0 & 0 & 1 & 0 & 0 & 1 & 0 & 0 & 0 & 0 \\
1 & 1 & 0 & 0 & 0 & 0 & 0 & 0 & 0 & 1 & 0 & 0 & 0 \\
0 & 1 & 1 & 0 & 0 & 0 & 0 & 0 & 0 & 0 & 1 & 0 & 0 \\
1 & 1 & 0 & 1 & 1 & 0 & 0 & 0 & 0 & 0 & 0 & 1 & 0 \\
0 & 1 & 1 & 0 & 1 & 1 & 0 & 0 & 0 & 0 & 0 & 0 & 1
\end{pmatrix}
\endgroup.
\end{equation}

\begin{definition}[Allowed programs and their length]\label{def:model}
A current wire value $X_i$ is represented by its linear form in the thirteen
inputs; initially $X_i=x_i$, so the matrix of current forms is $I_{13}$.
A directed CNOT gate $(c,d)$, with control $c$, target $d$, and $c\ne d$,
replaces $X_d$ by $X_d+X_c$ and leaves every other row unchanged. Thus every
ordered pair of distinct labelled coordinates is allowed: there are exactly
$13\cdot12=156$ gates, each of unit cost. A chronological list
$g_1,\ldots,g_L$ applies these row updates successively to $I_{13}$; its cost
is $L$. We use \emph{program} and \emph{word} interchangeably for such a
chronological gate list.
We define
\[
\cost(M)=\min\{L:\text{a legal length-$L$ program maps }I_{13}\text{ to }M\}.
\]
Equality with $M$ means equality of all thirteen output linear forms for all
$x\in\F^{13}$. There are no prepared zero inputs, additional coordinates,
Hadamard gates, measurements, or free input/output permutations. In
particular, the coordinate order is fixed throughout. We minimize the cost
of the full labelled operator, not its action on an encoding subspace.
\end{definition}
\prooffigure{cnot_semantics}{fig:cnot-semantics}{One allowed gate in all
three equivalent descriptions. The physical statement
$X_d\leftarrow X_d+X_c$ is the same as adding row $c$ to row $d$ over
$\F$. Only the target changes, so a length-$L$ word performs exactly $L$
row additions.}

\begin{theorem}\label{thm:main}
The matrix in \eqref{eq:matrix} has exact minimum length $\cost(M)=19$.
Equivalently, every legal program implementing $M$ has at least nineteen
gates, and a legal nineteen-gate program exists.
\end{theorem}
The proof is organized around the universal lower bound. The finite
exclusion needed after the structural reduction is reproducible from the
included source. Its rerun checks the same implementation, not a separately
implemented solver. Section~\ref{sec:formal} states the precise current
Lean-verification boundary.

A row $r$ is encoded by $\sum_{j=0}^{12}r_j2^j$. This fixes every matrix
entry and the bit order. For comparison with the machine-readable data,
the same matrix has row masks
\begin{equation}\label{eq:target}
(M_0,\ldots,M_{12})=(1,2,4,8,16,32,329,402,292,515,1030,2075,4150).
\end{equation}
Write $S=\{0,1,2,3,4,5\}$, let coordinate 8 be the \emph{central}
coordinate, and order the remaining coordinates as
$Y=(6,7,9,10,11,12)$. Use input names $(a,b,c,d,e,f)$ on $S$, $t$ on
coordinate 8, and $(y_0,\ldots,y_5)$ on $Y$. The operator is
\begin{align*}
(a',b',c',d',e',f')&=(a,b,c,d,e,f),& t'&=t+c+f,\\
y'_0&=y_0+a+d+t,&y'_1&=y_1+b+e+t,\\
y'_2&=y_2+a+b,&y'_3&=y_3+b+c,\\
y'_4&=y_4+a+b+d+e,&y'_5&=y_5+b+c+e+f.
\end{align*}
The coordinates in $S$ will be called \emph{sources}: their final rows are
unchanged basis rows. The coordinates in $Y$ will be called \emph{private
outputs}; ``private'' refers only to the displayed incidence pattern---each
row retains its own independent input $y_j$---and never means that $y_j$ is
zero or initialized specially. Let
\[
J=S\cup\{8\},\qquad
\pi:\F^{13}\longrightarrow\F^7
\]
be the coordinate projection that retains the coefficients on $J$, in the
order $(0,1,2,3,4,5,8)$, and deletes the six coefficients on $Y$. Thus an
initial private row has projection zero even though its physical input
$y_j$ is arbitrary.
\prooffigure{target_structure}{fig:target-structure}{Block structure of the
target transformation. A dot marks a retained coefficient equal to one.
Every private output also retains its own independent input $y_j$; the
projection $\pi$ merely omits that coefficient in the auxiliary seven-bit
problem. The picture restates the target equations and is not a circuit.}
\prooffigure{projection_split}{fig:projection-split}{The two views of one
physical row. Its full thirteen-bit value consists of seven retained
coefficients and six private-input coefficients. The projection $\pi$ keeps
the first block and forgets the second only for bookkeeping; it does not set
any physical input to zero and does not modify the circuit.}

For a program $\mathcal P$, define
\[
\begin{aligned}
u(\mathcal P)&=|\{s\in S:\text{some gate of $\mathcal P$ targets $s$}\}|,\\
v(\mathcal P)&=|\{y\in Y:\text{some gate of $\mathcal P$ uses $y$ as control}\}|.
\end{aligned}
\]
These count distinct labelled coordinates, not gate occurrences. We write
$u$ and $v$ when the program is understood.

\subsection{Notation guide}
The proof uses several representations of the same circuit. The following
tables are a lookup guide; every local symbol is also defined again where it
first enters an argument.
\begin{center}\begin{tabular}{lp{.71\linewidth}}\toprule
Symbol&Meaning\\\midrule
$x$&arbitrary physical input vector in $\F^{13}$\\
$X_i$&current linear form stored on labelled wire $i$\\
$M$&the fixed target matrix in~\eqref{eq:matrix}\\
$L$&length of the particular program currently under discussion\\
$\cost(M)$&minimum legal CNOT length of the target matrix\\
$S$&the six source coordinates $\{0,\ldots,5\}$\\
$t$&the central input, physically located on coordinate 8\\
$Y$&the six private-output coordinates $(6,7,9,10,11,12)$\\
$J$&the seven retained coordinates $S\cup\{8\}$\\
$\pi$&projection of a full row onto the coefficients indexed by $J$\\
$u$&number of distinct sources ever used as targets\\
$v$&number of distinct private outputs ever used as controls\\\bottomrule
\end{tabular}\end{center}

Later sections use $P$ for the seven initial projected basis forms,
$f_0,\ldots,f_6$ for the seven required projected increments, $F$ for their
ordered list, and $\alpha$ for free XOR complexity. The counting variables
$h$ and $\rho$ denote, respectively, the number of additional projected
forms and the number of unclassified repeated projected results. These names
are chosen to avoid confusion with the physical source variables
$(a,b,c,d,e,f)$.
\prooffigure{representation_layers}{fig:representation-layers}{The four
levels of description used in the proof. The physical word acts on thirteen
full rows. Projection keeps seven coefficients. A sector selects a smaller
core plus passive outputs. The finite search then stores only a sufficient
label for each core prefix. Every arrow is justified later; none changes the
physical synthesis problem.}

\subsection{How the lower bound works}
Figure~\ref{fig:map} separates the structural deductions from the remaining
finite exclusion. The proof has five stages:
\begin{enumerate}
\item the explicit witness gives the upper bound $\cost(M)\le19$;
\item projection by $\pi$, followed by forced-copy and restoration counts,
gives the two budgets $L\ge16+u$ and $L\ge16+v$;
\item these budgets and the separate $u=v=0$ argument place every
hypothetical program of length at most eighteen in one of 363 sectors;
\item catalogue, clean-carrier, passive-target, and remaining-cost tests
define a finite search in each sector; and
\item the C++ run finds no terminal state of total cost at most eighteen.
\end{enumerate}
Lean formalizes the general reduction and abstract search; the concrete
finite exclusion remains an external computation, as detailed in
Section~\ref{sec:formal}.
\prooffigure{proof_map}{fig:map}{Logical structure of the lower bound.
Every program of length at most eighteen must follow one of the displayed
branches; the mathematical reductions plus the reported exhaustive search
rule out every branch. The diagram is a roadmap, not a search certificate or
a claim that every displayed implication has been checked in Lean. The
nineteen-gate upper bound is separate.}

\section{Related results and provenance}
Patel, Markov, and Hayes prove the asymptotic worst-case CNOT count
$\Theta(n^2/\log n)$~\cite{pmh}. This does not determine the cost of
\eqref{eq:target} or an exact diameter at $n=13$.
Bu, Fan, and Joo give a general lower bound and prove that an $n$-cycle
permutation requires $3(n-1)$ CNOT gates~\cite{bfj}. In particular, a
13-cycle requires 36 gates, so the present nineteen-gate matrix is not a
diametral example. Their computable LMC estimate evaluates to 12 for our
matrix and for its transpose, inverse, and inverse transpose; it does not
directly establish the claimed value 19.

Sodhani and Parhi report a nineteen-gate circuit for an HGP $[[13,1]]$
encoder~\cite{sp}. Their benchmark table supplies an achieved count, not
an exclusion of all eighteen-gate realizations. Our local reconstruction
uses a length-three repetition-code parity-check matrix and fixed
elimination conventions. It is not author-supplied benchmark data.
Theorem~\ref{thm:main} concerns only \eqref{eq:target}; no stronger
benchmark-identity or priority claim is made. Standard XOR-DAG and
transposition arguments used below are not claimed as new general methods.

\section{An explicit upper bound}
Figure~\ref{fig:circuit} shows the witness in the original coordinate order.
It establishes $\cost(M)\le19$; the later sections establish
the central claim $\cost(M)\ge19$.
The following chronological program realizes $M$:
\begin{align*}
&(0,9),(1,9),(1,10),(2,10),(2,12),(5,12),\\
&(3,0),(0,11),(4,1),(1,11),(1,12),(1,7),(0,6),\\
&(8,7),(8,6),(2,8),(5,8),(4,1),(3,0).
\end{align*}
The successive new target values are
\begin{align*}
&a+y_2,\ a+b+y_2,\ b+y_3,\ b+c+y_3,\ c+y_5,\ c+f+y_5,\\
&a+d,\ a+d+y_4,\ b+e,\ a+b+d+e+y_4,\\
&b+c+e+f+y_5,\ b+e+y_1,\ a+d+y_0,\\
&b+e+t+y_1,\ a+d+t+y_0,\ c+t,\ c+f+t,\ b,\ a.
\end{align*}
These give precisely \eqref{eq:target}. Each gate is an involution, so the
target is invertible as well.
\prooffigure{cnot_circuit}{fig:circuit}{The chronological nineteen-gate
witness. The two final gates restore the temporarily changed source rows.
This diagram is generated from the same gate list as the Lean witness and
checked against all thirteen rows of $M$. It proves only the upper bound
$\cost(M)\le19$; it does not address optimality.}

\section{Projected XOR complexity}
Apply $\pi$ to each current row, discarding its six private-input
coefficients only in this auxiliary argument. No physical input is removed
or fixed: the original thirteen inputs remain arbitrary. The projected
problem is deliberately easier than the physical one: it remembers every
form once produced, permits unlimited reuse, and never overwrites a form.
Therefore the first appearances in any physical CNOT program form a valid
free XOR construction.
The initial nonzero projected forms are
$P=\{a,b,c,d,e,f,t\}$. The seven required new forms are
\begin{align*}
f_0=r_0&=a+d+t,&f_1=r_1&=b+e+t,&f_6=r_2&=c+f+t,\\
f_2=q_0&=a+b,&f_3=q_1&=b+c,&f_4=s_0&=r_0+r_1,
&f_5=s_1&=r_1+r_2.
\end{align*}
Thus $f_j$ is the required projected increment of private output $y_j$ for
$0\le j<6$, while $f_6$ is the required final projected value of the central
row. In the order $(f_0,\ldots,f_6)$ their bit masks are
\begin{equation}\label{eq:F}F=(73,82,3,6,27,54,100).\end{equation}
A free XOR program retains all previous forms and adds two available forms
at unit cost. Write $\alpha(F\mid P)$ for its minimum cost.

\begin{lemma}\label{lem:free} $\alpha(F\mid P)=10$.\end{lemma}
\begin{proof}
Let $H=\spanof(r_0,r_1,r_2)$. The coefficients of $d,e,f$ give
$\dim H=3$, and its masks are
$\{0,27,45,54,73,82,100,127\}$.
In $\F^7/H$, the classes of $a,b,c,t$ are independent, with
$d\equiv a+t$, $e\equiv b+t$, and $f\equiv c+t$.
The nine elements of $(P\cup F)\setminus H$ therefore have distinct
nonzero classes, and $P\cap H=\varnothing$.

The intuition is that a new form outside $H$ can create at most one new
independent difference inside its coset, whereas the targets inside $H$
require three independent directions. We now make that count precise.

Let $B_{\rm aux}$ be the distinct additional computed forms outside $H$ and
outside $P\cup F$. Let $W\subseteq H$ be spanned by differences of pairs in
$((P\cup F)\setminus H)\cup B_{\rm aux}$ belonging to the same coset.
Every computed form inside $H$ belongs to $W$: induct on the program.
Two parents inside $H$ stay in $W$; exactly one parent in $H$ cannot
produce an element of $H$; two parents outside $H$ whose sum is in $H$
give one of the defining differences. Thus $\dim W\ge3$.
Initially there are no such differences. Adding one element to a coset
increases the dimension of their span by at most one, because after
choosing one representative all new differences are generated by one
new difference and the old ones. Hence $|B_{\rm aux}|\ge3$.
The seven distinct target forms and at least three additional forms
require at least ten additions. Conversely, compute $a+d,b+e,c+f$,
then their sums with $t$, the two $q$ forms, and the two adjacent pair
sums of these three auxiliary forms. This uses ten additions.
\end{proof}

\prooffigure{quotient_argument}{fig:quotient}{The dimension obstruction
in Lemma~\ref{lem:free}. Equalities such as $d=a+t$ are identities of cosets
modulo $H$. The nine outside forms occupy distinct nonzero cosets;
each additional outside form contributes at most one new direction to $W$.
This proves the projected free-XOR lower bound, not by itself a lower bound
for full thirteen-wire CNOT programs.}

For a binary matrix $T$ whose columns represent independent input forms,
write $\alpha(T)$ for its free XOR complexity from the basis inputs. The next
lemma says that an XOR computation graph can be read backwards.
\begin{lemma}\label{lem:transpose}
For a binary $m\times n$ matrix $T$ without zero rows or columns,
free XOR complexity satisfies $\alpha(T^\top)=\alpha(T)+m-n$.
\end{lemma}
\begin{proof}
Here $n$ is the number of input forms and $m$ the number of requested output
forms. Remove unused computation nodes from a circuit with $s$ additions,
leaving $s'\le s$, and attach one pointer to each of the $m$ outputs. Every
remaining input
or computation node has at least one outgoing use. The total number of
uses, including multiple edges and output pointers, is $2s'+m$.
Reverse the edges. If a node $v$ had $d_v$ outgoing uses, it now has $d_v$
incoming values and sums them using $d_v-1$ additions. The total is
$\sum_v(d_v-1)=s'+m-n$. The parity of directed paths describes every
input-output coefficient; reversing edges transposes these coefficients.
This proves one inequality; applying it to $T^\top$ proves the reverse.
\end{proof}
The matrix with rows \eqref{eq:F} is square without zero rows or columns;
its transpose has rows
\begin{equation}\label{eq:FT}F^\top=(21,62,104,17,50,96,67).\end{equation}
Its free complexity is also ten. Transposing a full CNOT program reverses
its order and exchanges each control and target, preserving length.
The normalized coordinate orders used below are
$(0,1,2,3,4,5,8,6,7,9,10,11,12)$ and
$(6,7,9,10,11,12,8,0,1,2,3,4,5)$ respectively.
This is a length-preserving correspondence between two synthesis problems,
not a free operation available inside the original gate model.
\prooffigure{transpose_correspondence}{fig:transpose-correspondence}{How the
transpose reduction acts on a full CNOT word. Reversing the word and swapping
every control-target pair preserves its length and exchanges the activity
statistics $u$ and $v$ after the stated relabeling. It does not authorize a
free transpose or permutation within the original synthesis problem.}

\section{A budget for repeated forms}
Fix a full program and use the program statistics $u$ and $v$ defined above.
Every gate is classified by the projected value of its new target row. A
\emph{first appearance} creates a projected form not seen earlier. A
\emph{repeat} produces a projected form already seen. We write $h$ for the
number of additional first-appearing forms outside $P\cup F\cup\{0\}$, and
$\rho$ for repeated results that are neither a forced first copy nor a
credited source restoration.
\begin{lemma}\label{lem:budget}
Every length-$L$ realization satisfies $L\ge16+u$ and $L\ge16+v$.
More precisely, in one orientation,
\begin{equation}\label{eq:budget}
L=13+h+u+\rho,\qquad h\ge3,\quad \rho\ge0.
\end{equation}
\end{lemma}
\begin{proof}
Retain all previously observed projected forms, including initial zero.
First occurrences give a free XOR program with $7+h\ge10$ additions.
Each of the six private outputs starts with zero projection and ends
nonzero; its first nonzero update copies an already available projected
control, and is a repeat. Each changed source must finally restore its
original full row, producing another repeat. These $6+u$ events are
distinct. The remaining $\rho$ repeated events give~\eqref{eq:budget}.
Apply the same argument to the transposed word to obtain the bound in $v$.
\end{proof}
Thus $L\le18$ implies $u,v\le2$. If $u=2$, then
$L=18,h=3,\rho=0$. If $u=1$, then $h+\rho\le4$. The two lower bounds are
never added together.
\prooffigure{operation_budget}{fig:budget}{The exact accounting identity
and its zero-slack consequence. The six first copies are unavoidable
because the private outputs start with zero projection. This does not set
their independent physical inputs to zero. The figure explains the event
count behind Lemma~\ref{lem:budget}; the original and transposed bounds are
applied separately, not added.}

\begin{lemma}\label{lem:zero} If $u=v=0$, then $L\ge20$.\end{lemma}
\begin{proof}
Only sources can update the central row, and only sources or the central
row can update private outputs. Write the central row as $t+z$.
The five target-row groups counted below are disjoint, so their compulsory
costs may be added. Let $m_{be}$ count the coordinates in $\{b,e\}$ that are
ever present in $z$, and let $m_{ad}$ count those in $\{a,d\}$. Set
$m=m_{be}+m_{ad}$.
Each visited coordinate in these sets must be toggled on and off;
$c,f$ must each be toggled an odd number of times. Central updates cost
at least $2+2m$.
Outputs $y_0,y_1$ require at least two central transmissions in total
and direct additions for the $4-m$ unvisited required coordinates,
costing at least $6-m$. Outputs $y_2,y_3$ each cost at least two:
without central transmissions they require two distinct sources, and
with them their number of central transmissions is positive and even.
Output $y_4$ costs at least $\min(4,6-m)$, by separating the absence
of central transmissions from the case of at least two transmissions.
Similarly $y_5$ costs at least $4-m_{be}$.
These target groups are disjoint, so their costs sum to
$16+m_{ad}+\min(4,6-m_{be}-m_{ad})$.
For $m\le2$ this is $20+m_{ad}\ge20$; for $m\ge2$ it is
$22-m_{be}\ge20$.
\end{proof}
\prooffigure{zero_sector_cost}{fig:zero-sector}{The additive cost count in
Lemma~\ref{lem:zero}. The five boxes correspond to disjoint target-row
groups, so their lower bounds add. The final split into $m\le2$ and $m\ge2$
shows directly that the sector $u=v=0$ cannot contain a program of length at
most eighteen.}

\section{Covering sectors and auxiliary catalogues}
In normalized coordinates let $A\subseteq\{0,\ldots,5\}$ be the
sources that must actually change, and $R\subseteq\{0,\ldots,5\}$
the private outputs allowed to control. Members of $R$ need not actually
control. Put
\[
p=|A|,\qquad q=|R|.
\]
Thus $p$ is the number of selected changing sources and $q$ is the number of
private outputs placed in the searchable core. Every source outside $A$
stays fixed, and every output outside $R$ is passive. An \emph{orientation}
means either the original problem or its length-preserving relabelled
transpose.
\begin{definition}[Activity sector]
An activity sector is the data $(\text{orientation},A,R)$. It searches words
in which exactly the sources in $A$ are required to change at least once and
only private outputs in $R$ are permitted to act as controls. The set $R$ is
an allowance: some of its members may never control.
\end{definition}
All possibilities remaining after Lemmas~\ref{lem:budget}
and~\ref{lem:zero} are covered by the following families:
\begin{center}\begin{tabular}{llllr}\toprule
Actual case&Orientation&$|A|$&$|R|$&Sectors\\\midrule
$u=2$, $v\le2$&original&2&2&225\\
$v=2$, $u\le1$&transposed&2&1&90\\
$u=v=1$&original&1&1&36\\
$(u,v)=(1,0)$ or $(0,1)$&both&1&0&12\\\bottomrule
\end{tabular}\end{center}
Pad the set of actual controlling outputs to the displayed size as needed;
this enlarges a sector. Sources in $A$ must actually change, which is
checked at termination. The 363 sectors need not be disjoint.
\prooffigure{sector_cover}{fig:sectors}{The activity-pair coverage.
Padding the allowed set of controlling outputs enlarges a sector.
Transposition is used to cover activity patterns with two such outputs;
it is a length-preserving correspondence between problems. This proves only
coverage by sector parameters; it does not prove completeness of the later
normalized state search or the finite exclusion.}
\prooffigure{sector_anatomy}{fig:sector-anatomy}{Anatomy of one sector.
The searchable core contains the $p$ sources in $A$, the central row, and the
$q$ private outputs in $R$. The other $6-p$ sources are fixed controls; the
other $6-q$ outputs are passive targets and never control. The sector records
permission to control, not a claim that every member of $R$ actually does.}

\Needspace{13\baselineskip}
\paragraph{Sector and catalogue notation.}
\Needspace{12\baselineskip}
\begin{center}\begin{tabular}{lp{.70\linewidth}}\toprule
Symbol&Meaning\\\midrule
$A$&selected source indices that must actually be targeted\\
$R$&private-output indices permitted to act as controls\\
$p,q$&cardinalities $p=|A|$ and $q=|R|$\\
$\Omega$&all possible auxiliary projected forms outside $P\cup F\cup\{0\}$\\
$B_0$&one candidate auxiliary set tested by catalogue closure\\
$B_*$&the complete auxiliary set of a hypothetical short program\\
$\mathcal C_k$&all closure-valid auxiliary sets of cardinality $k$\\\bottomrule
\end{tabular}\end{center}

For the target list $F$ in either orientation, set
$\Omega=\F^7\setminus(P\cup F\cup\{0\})$, with $|\Omega|=113$.
Here $\Omega$ is the universe of possible nonzero auxiliary projected forms.
For $k=3,4$, enumerate every $k$-element subset $B_0\subseteq\Omega$.
Starting from $P$, repeatedly add an element of $F\cup B_0$ when it is
the sum of two available forms. Retain $B_0$ exactly when this closure
reaches all of $F\cup B_0$; call the resulting catalogue $\mathcal C_k$.
To see necessity for the short programs under consideration, take their
complete auxiliary set $B_*$. The budget gives $|B_*|=h\in\{3,4\}$. Order
$F\cup B_*$ by first appearance. At its first appearance,
each form is the sum of two forms already present in
$P\cup F\cup B_*$. The closure process can therefore reproduce this order
and reaches every element, so $B_*\in\mathcal C_{|B_*|}$. Conversely, a
closure sequence is itself a free XOR construction of all listed forms.
Monotonicity makes the final closure independent of the order in which
available additions are selected.
Equivalently, if $x_i$ denotes the $i$th form in first-appearance order,
there are earlier available parent forms $u_i$ and $v_i$ with
$x_i=u_i+v_i$.
The implementation examines all $\binom{113}{3}=234136$ triples and
$\binom{113}{4}=6438740$ quadruples. It obtains
\begin{center}\begin{tabular}{lrr}\toprule
Orientation&$|\mathcal C_3|$&$|\mathcal C_4|$\\\midrule
Original&298&21687\\Transposed&298&21823\\\bottomrule
\end{tabular}\end{center}
At a search prefix the observed auxiliary forms need only be contained
in some feasible full set. Capacity four permits subsets of both triples
and quadruples. These catalogues are regenerated from $F$ by the verifier;
their displayed sizes are diagnostics, not externally trusted inputs.
\prooffigure{catalogue_closure}{fig:catalogue-closure}{Why the catalogue test
is necessary for a short program. The complete projected trace orders every
required or auxiliary form by first appearance, proving that its full
auxiliary set belongs to $\mathcal C_3$ or $\mathcal C_4$. A search prefix
need only extend to such a set within its remaining capacity. The verifier
still enumerates every candidate triple and quadruple; the picture is not a
substitute for that enumeration.}

\section{The zero-slack clean-carrier argument}
This section applies only to $p=|A|=2$, where $h=3$ and $\rho=0$.
Write a full row as $(\phi,\mu)$, where $\phi\in\F^7$ is its projected part
and $\mu\in\F^6$ records its six private-input coefficients. Call the row
\emph{clean} when $\mu=0$, and \emph{dirty} otherwise. Thus clean does not
mean that $\phi$ is a basis form; it means only that no private coefficient
contaminates the row.
\begin{lemma}\label{lem:fresh}
Before a private output $j$ first acquires nonzero projection, its row
is its original private unit vector, and that private bit is absent
from every other row.
\end{lemma}
\begin{proof}
A zero-to-zero projected update would be an unbudgeted repeat, so this
output has not yet been targeted. Its first use as control would export
its fresh private bit while leaving the target projection unchanged.
That is neither a first nonzero copy nor a full source restoration, and
would again be an unbudgeted repeat.
\end{proof}
\begin{lemma}\label{lem:clean}
After the first occurrence of $\phi\notin P\cup\{0\}$, another occurrence
can only be a first nonzero copy to an uninitialized private output.
It adds that output's fresh private bit. Consequently a dirty first
occurrence cannot later produce a clean occurrence, and a clean occurrence
cannot be recreated after it disappears.
\end{lemma}
\begin{proof}
A source restoration gives a member of $P$, so the only admissible repeat
of $\phi$ is a first nonzero copy. Let $e_j\in\F^6$ be the private unit
vector belonging to output $j$. By Lemma~\ref{lem:fresh}, such a copy sends
$(\phi,\mu)$ to $(\phi,\mu+e_j)$ with the fresh bit $e_j$ absent from $\mu$.
Successive copies of this same projection cannot cancel those fresh bits.
To reach the final private row $(\phi,e_j)$, either $\phi$ is first created
at its final update on $j$, or a clean $\phi$ is copied into uninitialized
$j$. In particular, a first occurrence carrying $e_j$ on a different
row implies that $j$ was already initialized, forbidding the latter copy.
\end{proof}

The central target $f_6$ must therefore first occur on its own row in
its final full form and must remain there. A controlling private output
whose required projection has already occurred must either be final,
or still have zero projection with a clean carrier currently present
on a core row. For a passive output whose required projection has
occurred in the core history, that history must contain its clean form.
These are necessary conditions only.
Fixed sources cannot carry a clean nonunit form. A passive output cannot
carry a clean form at all: its own private bit cannot enter any controlling
row and hence cannot be cancelled there. If a target projection were
first created on a passive output, it could not subsequently return to
the core without an unbudgeted repeat. Removing passive updates therefore
does not invalidate the stated necessary conditions.
\prooffigure{clean_carrier}{fig:clean}{The fresh-private-bit mechanism
behind Lemmas~\ref{lem:fresh} and~\ref{lem:clean}. In the zero-slack case,
the only allowed repeat of a nonbasis projection is a first copy into an
uninitialized output, which adds a private bit that cannot cancel. These are
necessary pruning conditions only, not a sufficient characterization of
valid histories.}

\section{Exact elimination of passive targets}
Fix $p=|A|$ and $q=|R|$. A core history consists of updates to the $p$
changed sources, the central row, and the $q$ potentially controlling
outputs. Denote this gate sequence by $\gamma$ and its length by
$K=|\gamma|$. Its rows retain
$n_{\rm core}=7+q\le9$ bits, including the private bits
of those $q$ outputs. Together with the $6-p$ fixed sources, these rows
form an invertible $n_{\rm core}\times n_{\rm core}$ system at every step.
Private bits of
passive outputs cannot enter this system because those outputs never
control; their omission is justified by the sector alphabet.

Deleting passive-target gates is only a compression of the history; their
cost is charged back exactly. Let $D(\gamma)$ be the set of every distinct
retained row value that occurs on an allowable control during $\gamma$,
including the initial basis. Embed each seven-bit $f_j$ into
$\F^{n_{\rm core}}$ by appending zero private coefficients. Define
\[
\delta_D(f)=\min\{|B|:B\subseteq D,\ \bigoplus_{b\in B}b=f\},
\qquad E(D)=\sum_{j\notin R}\delta_D(f_j).
\]
Here $\delta_D(f)$ is the smallest number of historical controls whose XOR
is $f$, and $E(D)$ is the total minimum completion cost of all $6-q$ passive
outputs.
\Needspace{10\baselineskip}
\begin{center}\begin{tabular}{lp{.70\linewidth}}\toprule
Symbol&Meaning\\\midrule
$\gamma$&the ordered sequence of core-target gates after passive gates are erased\\
$K$&core length $|\gamma|$\\
$n_{\rm core}$&retained row width $7+q$\\
$D(\gamma)$&all distinct retained control values available during $\gamma$\\
$\delta_D(f)$&minimum subset size in $D$ whose XOR is $f$\\
$E(D)$&sum of these minimum costs over all passive outputs\\\bottomrule
\end{tabular}\end{center}
The initial basis ensures existence; repeated copies of one form cancel
and can be removed.
\begin{lemma}\label{lem:passive}
For a core history $\gamma$ ending at all prescribed core target rows,
the minimum cost of a full circuit extending that history is
$|\gamma|+E(D(\gamma))$.
\end{lemma}
\begin{proof}
Every update of a passive row adds a member of $D$, and its total added
value must be exactly $f_j$, including cancellation of any private
contamination. This proves the lower bound. Conversely, choose a minimum
subset for each $j$ and insert each transmission at a time when its
control value is available. Passive updates do not affect the core or
one another, so all insertions coexist and give the required full operator.
\end{proof}
The endpoint hypothesis in this lemma is essential.
\prooffigure{passive_elimination}{fig:passive}{A schematic passive-output
history. Dashed transmissions may be erased without changing the full core
rows, and inserted at the moments their control values exist. The displayed
three-term XOR is illustrative; the minimization uses every subset of the
actual historical control set $D$. The exact conclusion is conditional on a
fixed core history with the prescribed endpoint; the picture does not show
that all such core histories have been enumerated.}

The verifier computes $\delta_D$ by exact zero-one subset dynamic
programming. For a processed prefix of $D$, let $c(x)$ be the smallest
subset size with XOR $x$. Start with $c(0)=0$ and $c(x)=\infty$ otherwise;
for each $w\in D$ update
\[
c_{\rm new}(x)=\min(c_{\rm old}(x),1+c_{\rm old}(x+w)).
\]
Reading both old values in each pair $\{x,x+w\}$ prevents reuse of $w$.
Induction over $D$ proves the recurrence. There are at most $2^9=512$
values. Each of the $6-q$ passive outputs needs at least one operation,
so a candidate of cost at most eighteen has core length
\begin{equation}\label{eq:cap}K\le12+q.\end{equation}

\section{A remaining-cost bound}
At a search state, let $X$ be the current retained matrix and $Z$ its target.
For each row index $i$ with $X_i\ne Z_i$, the vector $X_i+Z_i$ is the
increment that must eventually be added to that row. Let $r$ be the number
of differing rows, and let $b$ be the number of distinct required increments
that occur as rows in neither endpoint matrix. These definitions are local
to the following lemma.
\begin{lemma}\label{lem:remaining}
For invertible binary matrices $X,Z$, let $r$ count differing rows and
let $b$ count distinct increments $X_i+Z_i$ for those rows that occur
as rows in neither $X$ nor $Z$. Every CNOT path between them has length
at least $r+\lceil b/2\rceil$.
\end{lemma}
\begin{proof}
Write its length as $r+k$. At most $k$ required rows can be updated
more than once. Removing their increments loses at most $k$ distinct
values. Each remaining increment must be available on a control at
the unique update of its target. Being absent from both endpoints,
it must have been created as an intermediate value. Reserve the last
update of each of the $r$ differing rows; these create only final rows.
All nonendpoint intermediate values are created among the other $k$
operations. Hence $b-k\le k$.
\end{proof}
Include fixed sources among the rows at both endpoints when applying
this to a core. Let $z_A$ be the number of chosen sources that have never
changed. Each such source still
equals its target but must change and return. Thus another bound is
$r+2z_A$. We use the maximum of the two valid bounds, not their sum:
\[B_{\rm rem}=\max\{r+2z_A,r+\lceil b/2\rceil\}.\]
The symbol $B_{\rm rem}$ therefore denotes a certified lower bound on the
number of future core gates still required.
\prooffigure{remaining_bound}{fig:remaining}{The counting argument for
Lemma~\ref{lem:remaining}. Here $b$ counts distinct required increments
absent from both endpoint row sets. The $k$ extra operations bound both
the discarded increments and the newly created intermediate values. This is
an admissible lower bound on the remaining cost, not a claim that the bound
equals the exact CNOT distance.}

\section{Finite search and completeness}
The verifier does not enumerate all $13\times13$ matrices or all length-18
gate words. For a fixed sector it enumerates only core histories, compresses
each prefix to a label containing all information relevant to the future,
and finally adds the exact passive completion cost.

For each sector, order core rows as chosen sources, central row, and
chosen controlling outputs. Choose fixed orderings
$A=(s_0,\ldots,s_{p-1})$ and $R=(j_0,\ldots,j_{q-1})$. Let $e_i$ denote the
$i$th standard basis vector of $\F^{n_{\rm core}}$. The initial core rows are
$e_{s_i}$, $e_6$, and $e_{7+\ell}$; their target rows are
$e_{s_i}$, $f_6$, and $e_{7+\ell}+f_{j_\ell}$, where
$0\le i<p$ and $0\le\ell<q$. Each $f_j$ is padded with $q$ trailing zeros.
Fixed controls are $e_s$ for $s\notin A$.
The retained coordinates are the seven projected coordinates followed by
the $q$ private bits belonging to $R$. Thus a full retained row is a
$n_{\rm core}=7+q$ bit mask. At every time, the core rows together with the
fixed source rows are the image of the initial $n_{\rm core}$ basis rows
under invertible row
operations; in particular they remain a basis and no genuine core update
can create the zero full row.

At each fixed core length store the complete label
\[\lambda=(\mathcal R,D,C,T,I,\rho_c),\]
where $\mathcal R$ is the tuple of current full core rows, $D$ the
historical full control set including the initial basis, $C$ the mask
of changed sources, $T$ the mask of credited full source restorations,
$I$ the mask of first nonzero private-output initializations, and $\rho_c$
the count of other projected repeats observed in the core.
The masks $C$ and $T$ have one bit for each of the $p$ selected sources;
$I$ has one bit for each of the $q$ selected private outputs.
Here $D$ contains every distinct full row value that has appeared in the
core history. Each such value was available as a control at the time it
appeared, although only current rows, never merely historical members of
$D$, may control a future core gate. The mathematical fields correspond
directly to the implementation:
\begin{center}\begin{tabular}{lp{.18\linewidth}p{.50\linewidth}}\toprule
Field&C++ member&Information retained\\\midrule
$\mathcal R$&\texttt{row}&current ordered core rows\\
$D$&\texttt{appeared}&all distinct full row values seen so far\\
$C$&\texttt{changed}&chosen sources targeted at least once\\
$T$&\texttt{returned}&one credited full restoration per chosen source\\
$I$&\texttt{copied}&one credited first nonzero initialization per core output\\
$\rho_c$&\texttt{extra}&all other repeated projected results in the core\\\bottomrule
\end{tabular}\end{center}
The sector, its orientation, $A$, $R$, and the core length are fixed outside
the label.
\prooffigure{search_state_anatomy}{fig:search-state}{What one search label
remembers. The current rows $\mathcal R$ determine legal next gates; the
historical set $D$ determines passive completion; $C$, $T$, $I$, and $\rho_c$
retain the finite budget information. Sector data $(A,R,p,q)$ and the layer
number $K$ are fixed outside the label. Equal labels at the same layer have
identical admissible futures.}

Enumerate every distinct control-target pair with target in the core
and control in the core or fixed sources. Reject a zero full new row.
Credit the first full restoration of a chosen source, and the first
nonzero projected initialization of a chosen private output. The choice
of first rather than last restoration is harmless: if one source returns
to its basis row $r_s$ times, exactly $r_s-1$ of those returns remain
additional repeats whichever return receives the credit. If a projected
result is in $P\cup\{0\}\cup\pi(D)$ and neither event is credited, increment
$\rho_c$ and require $\rho_c\le2-p$.
Here $\pi(D)=\{\pi(d):d\in D\}$ is the set of projections of historical
control values.

Let $B_{\rm obs}$ be the set of auxiliary projected forms observed in the
current core prefix. Require $B_{\rm obs}$ to extend to a feasible
catalogue set of size at most $5-p-\rho_c$. Update the rows and $D$; when
$p=2$, apply the clean-carrier conditions. Finally require the current
length plus $B_{\rm rem}$ not to exceed \eqref{eq:cap}.
Historical values in $D$ are never used as current core controls.

The direction of every rejection test is important. For a putative word of
length at most eighteen, the exact budget with $|A|=p$ gives
\begin{equation}\label{eq:shortbudget}
h+\rho\le5-p,\qquad \rho\le2-p,
\end{equation}
because $h\ge3$. Erasing passive updates can hide earlier projected
occurrences, so the core count $\rho_c$ is at most the global extra-repeat
count $\rho$. The auxiliary forms seen in the core are likewise a subset of
the global auxiliary set $B_*$. By the first-appearance argument above, $B_*$
belongs to a feasible catalogue, and
$|B_*|=h\le5-p-\rho\le5-p-\rho_c$. Hence both the repeat limit and the catalogue
extension test are relaxations of \eqref{eq:shortbudget}, not additional
restrictions on a true candidate. When $p=2$, Lemmas~\ref{lem:fresh}
and~\ref{lem:clean} make the clean-carrier tests necessary. Finally,
Lemma~\ref{lem:remaining} and the untouched-source bound are lower bounds on
the actual remaining core suffix, so pruning when the current length plus
$B_{\rm rem}$ exceeds \eqref{eq:cap} is also safe.

At every layer, inspect labels with all core rows final and every
chosen source marked changed. Such a label is a candidate exactly when
$K+E(D)\le18$. Equal labels at an equal layer can be merged because all
future transitions and tests depend only on the current rows, $D$, the four
counters or masks, and the fixed sector data. The terminal cost depends only
on $D$, the final rows, and the changed-source mask. Thus two equal labels at
the same layer have identical admissible continuations and terminal cost.
Hash collisions are resolved by complete label equality, so hashing does not
identify unequal labels.
\prooffigure{search_coverage}{fig:search-coverage}{The completeness pipeline
formalized in Lemma~\ref{lem:coverage}. A hypothetical short full word is
assigned to a sector, stripped to its core history, and followed through the
retained-label search. Every rejection test is one-sided, and passive gates
are charged back exactly by $E(D)$. Thus a true short realization would
produce a terminal candidate inspected by the verifier.}

\begin{lemma}\label{lem:coverage}
Any realization of $M$ with at most eighteen gates yields a terminal
candidate in at least one searched sector.
\end{lemma}
\begin{proof}
Let $w$ be such a realization.

\emph{Sector and orientation.}
The two budget bounds give $u,v\le2$, while Lemma~\ref{lem:zero} excludes
$u=v=0$. The activity-pair table therefore selects one of the 363 sectors,
possibly after the length-preserving transpose correspondence. Its set $A$
is exactly the sources targeted by the oriented word. Its set $R$ contains
every private output used as a control; padding $R$ only moves additional
outputs into the core and cannot invalidate the word.

\emph{Erasure and representation.}
Delete every gate whose target is outside the core. Such a passive coordinate
never controls, so these deletions leave the entire sequence of core rows
unchanged. Sources outside $A$ are never targeted and remain fixed controls.
Consequently every surviving gate has a core target and a control that is
either a current core row or a fixed source, exactly the alphabet enumerated
by the search. Passive private bits never enter a core row, so restricting
the surviving rows to the $n_{\rm core}=7+q$ retained coordinates loses no information
about this core history. If its length is $K$, the deleted part contains at
least one update for each of the $6-q$ passive outputs; hence
$K\le18-(6-q)=12+q$, as in \eqref{eq:cap}.

\emph{Induction over the core history.}
The initial search label is precisely the initial retained basis, with that
basis recorded in $D$ and all masks and counters zero. Suppose the label of a
genuine prefix survives at its layer. The next genuine gate is among the
enumerated control-target pairs. Its new full row is nonzero by invertibility.
Updating $\mathcal R$ and adjoining the new value to $D$ reproduces the next
genuine prefix. The definitions of $C,T,I$, and $\rho_c$ reproduce the credited
change, restoration, initialization, and repeat events; crediting the first
rather than the final restoration only changes which one of several returns
is exempted, not their total number. The repeat and catalogue tests are safe
by \eqref{eq:shortbudget} and the catalogue first-appearance argument. In the
zero-slack case the clean-carrier tests are necessary by
Lemmas~\ref{lem:fresh} and~\ref{lem:clean}. The remaining-cost test is no
larger than the genuine suffix. Thus the successor is retained, unless it is
merged with an equal label; label sufficiency then preserves the same
continuation. Induction reaches the complete erased core history.

\emph{Terminal cost.}
The final core rows are their prescribed full targets, and every member of
$A$ has been targeted, so the resulting label passes the terminal test. For
each passive output, the controls used by the deleted gates occur in $D$ and
their XOR is its required increment. Cancelling duplicate controls gives a
subset of $D$ no larger than the deleted list. Therefore
$E(D)$ is at most the number of deleted gates; by Lemma~\ref{lem:passive},
it is also the exact minimum passive completion cost. Hence
$K+E(D)\le |w|\le18$, and the label is a
terminal candidate.
\end{proof}

\paragraph{Finite exclusion.}
A fresh full run of the supplied implementation completes all sectors:
\begin{center}\begin{tabular}{lrrr}\toprule
Family&Sectors&Layer labels&Terminal labels\\\midrule
Two sources, two outputs&225&240830308&21684\\
Transposed: two sources, one output&90&3175146&1848\\
One source, one output&36&62738391&65679\\
One source, no controlling outputs&12&179149&1305\\\midrule
Total&363&306922994&90516\\\bottomrule
\end{tabular}\end{center}
No terminal label has total cost at most eighteen. These are sums of
layer sizes, not counts of distinct matrices or circuits. A sentinel
cost of 999 means no terminal was found, not a length-999 witness.
The successful exit is zero with the final marker
\texttt{VERIFIED\_FIXED\_MATRIX\_MINIMUM\_19}.
The implementation imposes no heuristic truncation of the searched layers.
An interrupted or failed run does not establish the exclusion.
Together with Lemma~\ref{lem:coverage}, this exhaustive result supplies the
lower bound. Its evidentiary form is source code plus a reported rerun, not an
independently checkable proof certificate.

\begin{proof}[Completion of Theorem~\ref{thm:main}]
A circuit of cost at most eighteen would contradict finite exclusion
by Lemma~\ref{lem:coverage}. The explicit nineteen-gate circuit supplies
the opposite inequality.
\end{proof}

\section{Reproducibility and formalization boundary}\label{sec:formal}
The proof separates a general mathematical reduction, its concrete C++
implementation, and the finite execution. Lean checks the general reduction
and derives the exact minimum from the explicit finite premise below.
The concrete tables and exclusion are supported by the external audits and
run evidence supplied under \texttt{anc/}; they are not replayed by Lean.
\prooffigure{verification_boundary}{fig:verification-boundary}{The updated
verification boundary. Lean checks the general reduction and the conditional
minimum theorem. Concrete catalogue completeness and finite exclusion are
supported externally; the C++ implementation is audited rather than
formally verified. No unconditional kernel proof is claimed.}

\subsection*{Finite C++ exclusion and implementation audit}
The source file \texttt{hgp13\_verify.cpp} regenerates all catalogues and
sectors and requires no data files or third-party C++ libraries.
Compile and run with
\begin{lstlisting}[numbers=none]
g++ -O3 -std=c++17 -pthread hgp13_verify.cpp -o hgp13_verify
./hgp13_verify 3
\end{lstlisting}
The rerun used GCC 11.4.0 in a Linux container with networking disabled.
Its source SHA-256 is
\begin{center}\scriptsize\ttfamily
2174f445df534455c30e9ae7b54b8dfff4fdbdce3ff4eaf6b7d7e557fa3f8e4c.
\end{center}
The accompanying evidence contains the full run log and source identity.
An additional full run compiled that exact source with the same GCC flags
in the pinned network-disabled container. The wrapper recorded the source
and binary SHA-256 values, command, exit code zero, and new transcript hash.
All 363 case records and family totals match the preserved run; the
manifest and audit are in
\texttt{anc/evidence/full\_provenance\_2026-09-27/}. This repeats the same
implementation rather than independently proving its semantics.
The reproducible script \path{anc/scripts/audit_full_run.py} checks all 363
case records, layer and family sums, two comparison transcripts, and equality
of the saved exported catalogue key sets with the literal Lean lookup trees.
An independent exhaustive enumerator, which does not include the verifier
source and adds all currently derivable forms in simultaneous rounds,
reproduces the same four key sets and catalogue counts. This is an
external cross-check of the concrete tables, not a kernel proof of their
completeness or of the C++ generator semantics.
Separate executable exports from C++ and Lean agree byte for byte on the
orientation, source and relay tuple of each of the 363 sectors.
Running the abstract Lean transition once for every sector yields 363
first-layer sizes that all agree with the saved C++ layer counts. A
source-hash-pinned instrumented C++ diagnostic and a separate Lean export
also agree on all 8,610 complete first-layer state records, including rows,
history bits, and counters. For sectors 234 and 359, source-hash-pinned
diagnostics exported every layer: respectively 6,668 and 2,208 complete
state records agree as sets with the candidate graphs whose transitions
are checked in Lean. These comparisons cover two sectors and do not
establish later-layer correspondence for the other 361 sectors.
The supplementary source audit records bounds for the packed arrays and
catalogue keys, the layer and 0/1-cost loop invariants, and the worker
join conditions. A separate UBSan diagnostic on sector 357 matched its
preserved full-run record; it does not certify the other sectors.
The pinned terminal-cost routine also agrees with a separate shortest-path
BFS oracle on 1,518 deterministic 7--9-bit histories and all 453,376
resulting cost entries; this bounded comparison is not a formal proof of
the C++ loop.
This agreement is repeatability evidence for the same implementation. The
log itself does not cryptographically attest which binary produced it, and
the search result is not a Lean kernel certificate.

\subsection*{Lean theorem and its finite premise}
The Lean development pins version 4.34.0 and uses its bundled standard
library. It specifies bit-vector rows, legal gates, full matrix execution,
the nineteen-gate witness, and the exact-minimum proposition. It also
develops the preservation of full core rows under erasure of passive-target
gates and exact row-update accounting. The kernel-checked development now
also proves free XOR complexity ten in both orientations, the necessity
of three distinct auxiliary projected values, full-word transposition,
and the unconditional bounds $L\geq16+u$ and $L\geq16+v$.
The exact budget $L=13+h+u+\rho$ and its short-word consequences are checked
in both orientations. For the concrete transpose free-XOR lemma, an
inductive first-appearance argument reduces the certificate to 198
ordered pairs of auxiliary values; the kernel checks both their
closed-set witnesses and coverage of all seven-bit pairs. This small
certificate does not certify the main CNOT exclusion. The general
rectangular free-XOR transposition principle is not yet formalized;
its concrete consequence here is proved independently.
The kernel also proves $L\geq20$ when $u=v=0$, by counting actual gates
and tracking coefficient parities. It verifies that the 363 sector parameter
records are distinct and cover every putative legal word of length at most
eighteen, preserving the chosen orientation's full operator and length.
The abstract auxiliary-catalogue lemma is checked: saturation is equivalent
to a valid first-appearance program, and every full realization supplies an
accepted auxiliary set in either orientation. The size bounds and abstract
subset test for partial histories are also checked. The concrete table
enumeration and its reported cardinalities remain outside this certificate.
At the saturated two-source budget, the kernel proves uniqueness of the
mandatory repeats and the fresh-private-bit invariant at every prefix.
It also proves that a clean nonbasis value must be globally new, and that
a repeated final nonbasis row is copied from its exact clean goal into an
intact private output. All three semantic core-history conditions are now
proved in both orientations: the seen central goal is already final; an
unfinished kept output has an intact row and a current exact clean carrier;
and a seen passive goal has its exact clean form in core history. Erasure
of passive-target gates preserves and reflects this full-row history.
The exact passive-cost theorem is also proved: the 0/1 recurrence computes
the minimum over distinct historical forms, all passive minima can be
attained simultaneously by legal insertions, and no completion with the
same core is shorter. Existence of the required representations follows
from the seven retained initial basis rows in both orientations.
For rows supported on a sector's retained coordinates,
\texttt{PackedCostRefinement.cost\_pack\_eq} proves that this exact
recurrence has the same value on full thirteen-bit rows and their packed
integer masks. The correspondence of the imperative C++ array update to
that recurrence is reviewed separately, not verified by Lean.
The theorem \texttt{SectorPackedCost.actual\_core\_cost\_pack\_eq}
instantiates this equality for the complete retained history of every
compatible full program and each of its passive target differences.
It also checks that those differences have no discarded-coordinate bits.
The stronger \texttt{SectorPackedCost.indexed\_history\_cost\_pack\_eq}
applies to every indexed sector path, including arbitrary terminal paths
quantified over in the external finite premise. Its complete full-row
history is supported, and its packed 0/1 cost equals the full-row cost.
\texttt{CostSetInvariant.cost\_eq\_of\_mem\_iff} proves that exact cost
depends only on which historical forms occurred. Its packed corollary
shows that deleting repeated supported forms leaves the cost unchanged,
matching the mathematical role of the C++ history bitset.
For every indexed sector path, the theorem
\texttt{IndexedHistoryBits.indexed\_history\_bit\_iff} identifies each
bit in the abstract search history exactly with membership of the
corresponding supported full row in the complete retained history.
Define $\mathcal H_s(h)$ as the supported full rows whose packed indices
are set in history bitset $h$. The theorem
\texttt{StateHistoryCost.indexed\_cost\_eq\_decoded} proves, for every
indexed path and every passive target, that its full historical 0/1 cost
equals the cost computed from $\mathcal H_s(h)$ at the terminal state.
Thus the mathematical terminal cost depends only on the complete state.
The imperative C++ loop remains an external implementation of that
packed recurrence.
The two remaining-cost bounds and their maximum are proved for full
13-coordinate states, with sound pruning and an application to actual
core suffixes. Lean now also proves that any legal short full word maps to
an accepted path of the abstract search when the concrete prefix lookup is
complete. For every sector, the exact passive cost of that path's full-row
history is bounded by the length of the original word. The theorem
\texttt{ConcreteFiniteBoundary.exact\_minimum\_of\_external\_finite\_result}
therefore concludes the minimum nineteen from one explicitly named premise,
\texttt{ExternalFiniteResult}. Here $P_s$ denotes the passive-output indices of sector $s$, and $M_s$
its oriented target matrix. The premise asserts complete concrete prefix
lookups and, for every sector $s$ and depth $k$ within its cap, that each terminal
state $z$ in the deduplicated abstract frontier $F_{s,k}$ satisfies
\[
  k+\sum_{j\in P_s}\operatorname{cost}_{\mathcal H_s(z.\mathrm{history})}
       (e_j\oplus M_s[j])>18.
\]
\texttt{FiniteFrontierBoundary.no\_cheap\_terminal\_of\_frontier}
transfers this state-and-depth condition to the path-quantified premise
of the general lower-bound theorem. The unconditional proposition is
\begin{lstlisting}[numbers=none]
def NoShortProgram : Prop :=
  forall p : Program, Legal p -> p.length <= 18 ->
    Not (run identity p = target)
\end{lstlisting}
At this draft stage it is an open Lean proposition, not an added axiom.
The conditional theorem does not make the external finite premise
kernel-verified. Its data and implementation correspondence are reviewed
in \texttt{anc/evidence/cpp\_refinement\_audit.md}; a static review is not a
formal proof of C++ semantics. The build and transitive-axiom audit record
which mathematical statements Lean accepts. Full kernel replay of all
sectors is outside this project's scope. No \texttt{sorry} or native
reduction trust axiom occurs in the accepted development.

\section*{Acknowledgements}
The author acknowledges OpenAI, the Astra AI assistant, and the CayleyPy
project. AI assistance was used in preparing the manuscript, examining the
computational evidence, and developing the Lean formalization with its explicit external finite premise.
Responsibility for the final mathematical claims and their presentation
rests with the author.

\begin{thebibliography}{9}
\bibitem{pmh} K. N. Patel, I. L. Markov, and J. P. Hayes,
``Optimal synthesis of linear reversible circuits,''
\emph{Quantum Information and Computation} 8 (2008), 282--294.
\url{https://www.rintonpress.com/xxqic8/qic-8-34/0282-0294.pdf}.
\bibitem{bfj} A. Bu, E. Fan, and R. Joo,
``Minimum Synthesis Cost of CNOT Circuits,'' arXiv:2408.07898 (2024).
\url{https://arxiv.org/abs/2408.07898}.
\bibitem{sp} A. Sodhani and K. K. Parhi,
``Noise-Aware Synthesis of Quantum LDPC Encoder Circuits via Two-Sided
Hamming Descent,'' arXiv:2607.04462 (2026).
\url{https://arxiv.org/abs/2607.04462}.
\end{thebibliography}

\appendix
\section{Complete source bundle of the current Lean development}\label{app:lean}
The complete Lean sources are supplied under \texttt{anc/lean/}, with
the pinned toolchain and Lake configuration. The package also includes
\texttt{anc/verify\_lean.py}, source-bound build and axiom-audit evidence,
the C++ run records, audit scripts, and \texttt{anc/lean-project.zip}.
No compiled cache is needed. The ZIP is a separate source-package file;
compiling this Prism project does not embed it as a PDF attachment.
The table lists the principal modules; the source archive includes their
dependencies and the full development.
{\footnotesize\begin{longtable}{@{}p{.40\linewidth}p{.56\linewidth}@{}}\toprule
Module&Role\\\midrule\endhead
\texttt{HGP13.lean}&Full matrix semantics, exact witness, and the target
minimum-length proposition.\\
\texttt{Passive.lean}&Preservation of full core rows and exact row-update
accounting.\\
\texttt{Certificate.lean}&A sound interface for complete finite layer
certificates.\\
\texttt{FreeXor.lean}&Exact free XOR complexity ten and three necessary
auxiliary values for the original masks.\\
\texttt{Projection.lean}&Transfer from arbitrary full-matrix programs to
the seven-bit free XOR model.\\
\texttt{CountHelpers.lean}&List counting and duplicate-removal lemmas.\\
\texttt{Budget.lean}&First-appearance and repeat accounting; the bound
$L\geq16+u$.\\
\texttt{ExactBudget.lean}&Exact operation budget and short-word saturation.\\
\texttt{Transposition.lean}&Coefficientwise transposition and relabeling
of arbitrary full CNOT words.\\
\texttt{TransposeFree.lean}&First-appearance restriction and soundness of
closed-set certificates.\\
\texttt{TransposeCertificate.lean}&Kernel-checked finite certificates and
complete coverage of auxiliary pairs.\\
\texttt{TransposeTarget.lean}&Concrete transposed target and the parameter
correspondence.\\
\texttt{TransposeBudget.lean}&Concrete transpose XOR minimum, the bound
$L\geq16+v$, and exact budgets.\\
\texttt{ZeroSector.lean}&Coefficient parities and the analytic lower
bound twenty for $u=v=0$.\\
\texttt{SectorCover.lean}&Distinctness and complete coverage of the 363
sector parameter records, with full-word semantics.\\
\texttt{Catalogues.lean}&Abstract saturation completeness, actual auxiliary
sets, size budgets, and the partial-history subset test.\\
\texttt{TightRepeats.lean}&Saturated repeat classification, fresh private
bits, and local restrictions on clean copies.\\
\texttt{CleanPersistence.lean}&Current clean carriers, all three semantic
core-history conditions, and exact history preservation under erasure.\\
\texttt{PassiveCost.lean}&Exact 0/1 XOR costs, pair cancellation, and total
passive-cost lower bounds.\\
\texttt{PackedCostRefinement.lean}&Equality of full-row and packed-mask
passive-cost recurrences for supported rows.\\
\texttt{SectorPackedCost.lean}&Supported retained histories and passive
target differences, with exact packed costs for arbitrary indexed paths.\\
\texttt{CostSetInvariant.lean}&Exact cost depends only on historical set
membership; duplicate packed forms can be removed.\\
\texttt{IndexedHistoryBits.lean}&Exact correspondence between abstract
history bits and full historical rows for every indexed path.\\
\texttt{DepthFrontier.lean}&Exact depth-dependent breadth-first reachability
after complete-state deduplication and inclusion of every accepted path.\\
\texttt{StateHistoryCost.lean}&Exact reconstruction of full historical
passive cost from the abstract terminal state's history bits.\\
\texttt{FiniteFrontierBoundary.lean}&State-and-layer exclusion implies
the path-quantified finite premise of the lower-bound theorem.\\
\texttt{PassiveInsertion.lean}&Legal historical insertions and the attained
minimum completion cost for each fixed core.\\
\texttt{RemainingBound.lean}&Both remaining-cost bounds, their sound
maximum, and application to actual core suffixes.\\
\texttt{FullWordPath.lean}&A short full word follows an accepted abstract
search path, conditional on concrete prefix completeness.\\
\texttt{AbstractFiniteExclusion.lean}&General terminal-cost transfer and
the minimum-length implication from explicit finite facts.\\
\texttt{ConcreteFiniteBoundary.lean}&The conditional theorem for the two
literal original and transposed lookup tables.\\
\texttt{Audit.lean}&Transitive axiom checks for the accepted declarations.\\\bottomrule
\end{longtable}}
``Complete source bundle'' describes file coverage. The central lower-bound
proposition remains open without the external finite premise, as stated in
Section~\ref{sec:formal}; attaching all current sources does not turn the
external computation into a kernel proof.

\section{Full finite verifier}
The following is the exact rerun source, not pseudocode.
\lstinputlisting[language=C++,basicstyle=\ttfamily\fontsize{8}{8.8}\selectfont]{hgp13_verify.cpp}
\end{document}
